% ============================================================================
%  Full mathematical breakdown of the ECR (Extended Control Regions) method,
%  specialised to the Lorenz system observed on a Poincare surface of section.
%
%  Written to follow on from the introductory paragraphs of
%  \section{ECR Data-Driven Targeting}\label{sec:ecr-method}
%
%  Requires: amsmath, amssymb, bm (optional), algorithm/algorithmic (optional).
%  Numerical values quoted here were produced by the accompanying MATLAB
%  implementation (ECR_III_MATLAB), integrator RK4 with dt = 10^{-3} and
%  Newton-refined section crossings.
%
%  NOTATION WARNING.  The Lorenz system has a coordinate called z and the ECR
%  formalism has an observed state called z.  Throughout this section the
%  observed (controller) state is set in bold, \mathbf{z}_n, while x, y, z
%  without decoration are the three Lorenz coordinates.  If you prefer to
%  remove the clash entirely, rename the observed state to \bm{\zeta}_n.
% ============================================================================

\subsection{Plant, surface of section and the controlled map}
\label{subsec:ecr-plant}

The Lorenz system is written in full, since the second equation is what defines
the surface of section used below,
%
\begin{align}
\dot{x} &= \sigma\,(y-x), \label{eq:lorenz-x}\\
\dot{y} &= \rho\,x - y - x z, \label{eq:lorenz-y}\\
\dot{z} &= x y - \beta z, \label{eq:lorenz-z}
\end{align}
%
with state $\mathbf{u}=[x\;\,y\;\,z]^{\mathsf T}\in\mathbb{R}^{3}$ and control
parameter vector
%
\begin{equation}
\mathbf{p}=[\sigma\;\,\rho\;\,\beta]^{\mathsf T}\in\mathbb{R}^{r},
\qquad r=3 .
\label{eq:param-vector}
\end{equation}
%
Let $\varphi_t(\mathbf{u};\mathbf{p})$ denote the flow of
\eqref{eq:lorenz-x}--\eqref{eq:lorenz-z}. The surface of section is the
half-plane
%
\begin{equation}
\Sigma=\bigl\{\mathbf{u}\in\mathbb{R}^{3}\;:\;y=y_{\mathrm{PSS}},\;\;\dot y<0\bigr\},
\qquad
y_{\mathrm{PSS}}=\sqrt{\beta_{\mathrm{nom}}(\rho_{\mathrm{nom}}-1)}=8.4853 ,
\label{eq:section}
\end{equation}
%
i.e.\ the $y$-coordinate of the non-trivial equilibria $C^{\pm}$. The sign
condition selects one of the two transversal branches; by \eqref{eq:lorenz-y}
transversality on $\Sigma$ requires
%
\begin{equation}
\dot y\big|_{\Sigma}=\rho x - y_{\mathrm{PSS}} - x z \neq 0 ,
\label{eq:transversality}
\end{equation}
%
which holds on the whole attracting set of interest. The branch $\dot y<0$ is
the one that reproduces the target and the state RMS quoted in the source
paper; on the branch $\dot y>0$ the equilibrium $C^{+}$ itself lies in $\Sigma$
and the fixed-point problem below degenerates onto it.

With the first return time
%
\begin{equation}
\tau(\mathbf{u};\mathbf{p})=\min\bigl\{t>0:\varphi_t(\mathbf{u};\mathbf{p})\in\Sigma\bigr\},
\label{eq:return-time}
\end{equation}
%
the Poincar\'e map is $\mathbf{u}\mapsto\varphi_{\tau}(\mathbf{u};\mathbf{p})$.
Because $y\equiv y_{\mathrm{PSS}}$ on $\Sigma$, the section state is fully
described by two coordinates, and the \emph{observed state} of the controller
is taken as
%
\begin{equation}
\mathbf{z}_n=\bigl[x(t_n)\;\;z(t_n)\bigr]^{\mathsf T}\in\mathbb{R}^{N},
\qquad N=2,
\label{eq:observed-state}
\end{equation}
%
where $t_n$ is the instant of the $n$-th crossing. The controlled dynamics
therefore reduce to the discrete map
%
\begin{equation}
\mathbf{z}_{n+1}=\mathbf{G}(\mathbf{z}_n,\mathbf{p}_n),
\label{eq:map}
\end{equation}
%
in which the parameter vector is held constant over the whole flight
$t\in[t_n,t_{n+1})$ and may be changed only at crossings. Equation
\eqref{eq:map} is the object the ECR method learns; no closed form for
$\mathbf{G}$ is ever constructed.

\paragraph{Delay-coordinate variant.}
When only a scalar measurement $s(t)=g(\mathbf{u}(t))$ is available, the
observed state is the delay embedding evaluated at the crossings,
%
\begin{equation}
\mathbf{z}_n=\bigl[s(t_n)\;\;s(t_n-T)\;\;\cdots\;\;s\bigl(t_n-(N-1)T\bigr)\bigr]^{\mathsf T},
\label{eq:delay-embedding}
\end{equation}
%
with delay $T$ and embedding dimension $N$. By Takens' theorem
\eqref{eq:delay-embedding} is generically diffeomorphic to the section state,
so \eqref{eq:map} still holds --- with the caveat discussed in
Remark~\ref{rem:delay}.

\subsection{Admissible parameter set}
\label{subsec:ecr-pi}

Each parameter may be perturbed about its nominal value within a prescribed
range,
%
\begin{equation}
\Pi=\Bigl\{\mathbf{p}\in\mathbb{R}^{r}\;:\;
\bigl|p^{i}-p^{i}_{\mathrm{nom}}\bigr|<\delta p^{i}_{\max},\;\;i=1,\dots,r\Bigr\},
\label{eq:admissible-set}
\end{equation}
%
with, for the Lorenz system,
$\mathbf{p}_{\mathrm{nom}}=[10\;\;28\;\;8/3]^{\mathsf T}$ and
$\delta\mathbf{p}_{\max}=[0.30\;\;0.84\;\;0.08]^{\mathsf T}$, i.e.\ $3\%$ of each
nominal value.  That budget is a design choice carried over from the source
paper, not something the method dictates: $\Pi$ need only be small enough to
leave the attractor intact, and both the extent of the control regions and the
reaching time scale with how large it is.  It is convenient to work with the
normalised perturbation
%
\begin{equation}
\Delta\mathbf{p}_n=\bigl(\mathbf{p}_n-\mathbf{p}_{\mathrm{nom}}\bigr)\oslash\delta\mathbf{p}_{\max},
\qquad
\mathbf{p}_n\in\Pi \iff \|\Delta\mathbf{p}_n\|_{\infty}\le 1 ,
\label{eq:normalised-perturbation}
\end{equation}
%
where $\oslash$ is the elementwise quotient. Every learned model below outputs
$\Delta\mathbf{p}$, so admissibility is a box test on the model output.

\subsection{Target orbit and its local geometry}
\label{subsec:ecr-target}

The target is a fixed point of the controlled map at nominal parameters,
%
\begin{equation}
\mathbf{z}^{*}=\mathbf{G}(\mathbf{z}^{*},\mathbf{p}_{\mathrm{nom}}),
\label{eq:fixed-point}
\end{equation}
%
which corresponds to a period-1 unstable periodic orbit of the flow (not to an
equilibrium: the equilibria of \eqref{eq:lorenz-x}--\eqref{eq:lorenz-z} do not
lie on the branch $\dot y<0$ of $\Sigma$). It is located by Newton iteration on
the residual $\mathbf{F}(\mathbf{z})=\mathbf{G}(\mathbf{z},\mathbf{p}_{\mathrm{nom}})-\mathbf{z}$,
%
\begin{equation}
\mathbf{z}^{(k+1)}=\mathbf{z}^{(k)}-\bigl(\mathbf{J}-\mathbf{I}\bigr)^{-1}\mathbf{F}\bigl(\mathbf{z}^{(k)}\bigr),
\qquad
\mathbf{J}=\frac{\partial\mathbf{G}}{\partial\mathbf{z}}\bigg|_{(\mathbf{z}^{*},\,\mathbf{p}_{\mathrm{nom}})},
\label{eq:newton}
\end{equation}
%
the Jacobian being evaluated by central differences of the numerically
integrated map. This yields
%
\begin{equation}
\mathbf{z}^{*}=\begin{bmatrix}14.2522\\ 39.7867\end{bmatrix},
\qquad
\tau^{*}=1.5587,
\qquad
\mathbf{J}=\begin{bmatrix}0.9594 & 1.7943\\ 2.0070 & 3.7536\end{bmatrix},
\label{eq:target-values}
\end{equation}
%
with eigenvalues
%
\begin{equation}
\lambda_1=4.7129,
\qquad
\lambda_2\approx 1.3\times10^{-9},
\label{eq:eigenvalues}
\end{equation}
%
so the target is a saddle whose stable direction is contracted so violently
that the return map is effectively one-dimensional --- the reason why control
regions on this section appear as thin filaments along the attractor.

The actuation available in one step is described by the sensitivity matrix
%
\begin{equation}
\mathbf{B}=\frac{\partial\mathbf{G}}{\partial\mathbf{p}}\bigg|_{(\mathbf{z}^{*},\,\mathbf{p}_{\mathrm{nom}})}
=\begin{bmatrix}
0.1739 & -2.3023 & -4.0222\\
0.2353 & -3.9651 & -10.3387
\end{bmatrix},
\qquad \operatorname{rank}\mathbf{B}=2 ,
\label{eq:sensitivity}
\end{equation}
%
whose singular values are $11.99$ and $0.659$. Since
$\operatorname{rank}\mathbf{B}=N$, the one-step image can be displaced in any
direction of the section plane; the reachable displacement under the full
parameter budget,
%
\begin{equation}
\bigl|\mathbf{B}\bigr|\,\delta\mathbf{p}_{\max}=[2.31\;\;4.23]^{\mathsf T},
\label{eq:reachable-displacement}
\end{equation}
%
is an order of magnitude larger than the target radius $\delta=0.30$ used
below, which is what makes single-step targeting into $T_0$ possible at all.
Note that $\sigma$ is by far the weakest actuator: alone it moves the image by
less than $\delta$.

\subsection{Control regions}
\label{subsec:ecr-regions}

Let $\delta>0$ be the target radius and define the $\delta$-ball
$\mathcal{B}_\delta=\{\mathbf{z}:\|\mathbf{z}-\mathbf{z}^{*}\|<\delta\}$.

\begin{definition}[Target region, $T_0$]
\label{def:T0}
\begin{equation}
T_0=S_0=\Bigl\{\mathbf{z}\in\mathcal{B}_\delta\;:\;
\exists\,\mathbf{p}\in\Pi \;\text{such that}\; \mathbf{G}(\mathbf{z},\mathbf{p})\in\mathcal{B}_\delta\Bigr\}.
\label{eq:def-T0}
\end{equation}
$T_0$ is thus the set of states that are both close to the target and can be
\emph{kept} close by an admissible perturbation; it is the region in which the
local (fine) controller operates.
\end{definition}

\begin{definition}[Extended control regions, $T_i$]
\label{def:Ti}
For $i=1,2,\dots,K$,
\begin{equation}
T_i=\Bigl\{\mathbf{z}\notin\textstyle\bigcup_{j<i}T_j\;:\;
\exists\,\mathbf{p}\in\Pi \;\text{such that}\; \mathbf{G}(\mathbf{z},\mathbf{p})\in T_{i-1}\Bigr\}.
\label{eq:def-Ti}
\end{equation}
Equivalently, $T_i$ is the set of states that are one admissible step from
$T_{i-1}$ and not already members of a lower region. The regions are disjoint
by construction, and
\begin{equation}
\mathbf{z}_n\in T_i \;\Longrightarrow\;
\exists\,\{\mathbf{p}_n,\dots,\mathbf{p}_{n+i-1}\}\subset\Pi:\;
\mathbf{G}\bigl(\cdots\mathbf{G}(\mathbf{z}_n,\mathbf{p}_n)\cdots,\mathbf{p}_{n+i-1}\bigr)\in T_0 ,
\label{eq:reachability}
\end{equation}
i.e.\ membership of $T_i$ certifies capture of the target region in at most $i$
steps. The union $\bigcup_{i=0}^{K}T_i$ is the \emph{activation region} of the
controller: outside it the nominal parameters are applied.
\end{definition}

\begin{remark}
Definitions~\ref{def:T0} and \ref{def:Ti} are stated with $\mathbf{G}$, but they
are \emph{existence} statements about admissible perturbations, not formulae.
Section~\ref{subsec:ecr-data} replaces each existential quantifier by evidence
drawn from measured transitions, which is what makes the method model-free.
\end{remark}

\subsection{Data model and the empirical construction of the regions}
\label{subsec:ecr-data}

The experiment is a random-perturbation sweep: the plant is run from many
initial conditions on the attractor while, at every crossing, a parameter
vector is drawn independently and uniformly from $\Pi$. The resulting data set
is a collection of transition triples
%
\begin{equation}
\mathcal{D}=\bigl\{(\mathbf{z}_n,\mathbf{p}_n,\mathbf{z}_{n+1})\bigr\}_{n=1}^{M},
\qquad \mathbf{p}_n\sim\mathcal{U}(\Pi),
\qquad \mathbf{z}_{n+1}=\mathbf{G}(\mathbf{z}_n,\mathbf{p}_n).
\label{eq:dataset}
\end{equation}
%
Region membership is then decided on $\mathcal{D}$ alone. Level $0$ uses
Definition~\ref{def:T0} directly, the applied $\mathbf{p}_n$ serving as the
witness of the existential quantifier:
%
\begin{equation}
\mathcal{I}_0=\bigl\{n\;:\;\|\mathbf{z}_n-\mathbf{z}^{*}\|<\delta
\;\wedge\;\|\mathbf{z}_{n+1}-\mathbf{z}^{*}\|<\delta\bigr\}.
\label{eq:index-T0}
\end{equation}
%
For $i\ge 1$ the test requires membership of the \emph{successor} in $T_{i-1}$,
a set that is itself known only through data. This is resolved by using the
already-trained level-$(i-1)$ model as the membership oracle. Writing
$\mathcal{M}_{i-1}:\mathbb{R}^{N}\to\{0,1\}$ for that oracle (defined in
\eqref{eq:membership-oracle} below),
%
\begin{equation}
\mathcal{I}_i=\Bigl\{n\;:\;\mathcal{M}_{i-1}(\mathbf{z}_{n+1})=1
\;\wedge\;\mathcal{M}_{j}(\mathbf{z}_{n})=0 \;\;\forall j<i \Bigr\}.
\label{eq:index-Ti}
\end{equation}
%
The training pairs of region $i$ are $\{(\mathbf{z}_n,\Delta\mathbf{p}_n)\}_{n\in\mathcal{I}_i}$:
the state, and the normalised perturbation that was observed to move it into
the next lower region. Construction proceeds bottom-up, $i=0,1,\dots,K$, each
level supplying the oracle for the next.

\subsection{Clustering of a control region (ECR-III)}
\label{subsec:ecr-clustering}

ECR-I models the whole activation region with a single network and ECR-II one
network per region $T_i$. ECR-III first partitions each region into clusters,
on the grounds that data belonging to the higher regions is scattered and a
single interpolant across disjoint patches is meaningless.

Let $\mathcal{Z}_i=\{\mathbf{z}_n\}_{n\in\mathcal{I}_i}$. Choose a radius $r_i$
and define on $\mathcal{Z}_i$ the relation
%
\begin{equation}
\mathbf{z}\sim\mathbf{z}' \iff \|\mathbf{z}-\mathbf{z}'\|\le r_i .
\label{eq:radius-relation}
\end{equation}
%
The clusters $C_{i1},\dots,C_{in_i}$ are the equivalence classes of the
transitive closure of \eqref{eq:radius-relation}; equivalently, two points lie
in the same cluster iff they are joined by a chain of points spaced by at most
$r_i$ (single-linkage at threshold $r_i$). The radius is selected from the
empirical distribution of pairwise distances: with
%
\begin{equation}
h_i(\cdot)=\text{histogram of }\bigl\{\|\mathbf{z}-\mathbf{z}'\|:\mathbf{z},\mathbf{z}'\in\mathcal{Z}_i,\;\mathbf{z}\neq\mathbf{z}'\bigr\},
\label{eq:distance-histogram}
\end{equation}
%
$r_i$ is taken at the first local minimum of $h_i$ following its first local
maximum --- the gap that separates intra-cluster from inter-cluster distances.

Each cluster is summarised by its sample mean and covariance,
%
\begin{equation}
\bm{\mu}_{ij}=\frac{1}{|C_{ij}|}\sum_{\mathbf{z}\in C_{ij}}\mathbf{z},
\qquad
\bm{\Sigma}_{ij}=\frac{1}{|C_{ij}|-1}\sum_{\mathbf{z}\in C_{ij}}
(\mathbf{z}-\bm{\mu}_{ij})(\mathbf{z}-\bm{\mu}_{ij})^{\mathsf T}
+\varepsilon\,\frac{\operatorname{tr}\bm{\Sigma}_{ij}}{N}\mathbf{I},
\label{eq:cluster-stats}
\end{equation}
%
the last term being a ridge that keeps $\bm{\Sigma}_{ij}$ invertible for the
thin, nearly one-dimensional clusters produced by a strongly contracting map.
The region of a current state is then decided by the \emph{normalised
Mahalanobis distance}
%
\begin{equation}
\mathrm{NMD}^{2}_{ij}(\mathbf{z}_n)=
(\mathbf{z}_n-\bm{\mu}_{ij})^{\mathsf T}\,\underline{\mathbf{N}}_{ij}\,(\mathbf{z}_n-\bm{\mu}_{ij}),
\label{eq:nmd}
\end{equation}
%
where $\underline{\mathbf{N}}_{ij}$ is the normalised covariance matrix. The
normalisation is not fixed by the source paper; the implementation supports
%
\begin{equation}
\underline{\mathbf{N}}_{ij}=
\begin{cases}
\bigl(\bm{\Sigma}_{ij}/\det(\bm{\Sigma}_{ij})^{1/N}\bigr)^{-1} & \text{unit determinant (default)},\\[2pt]
\bm{\Sigma}_{ij}^{-1} & \text{plain Mahalanobis},\\[2pt]
\operatorname{diag}\bigl(1/\Sigma_{ij,11},\dots,1/\Sigma_{ij,NN}\bigr) & \text{normalised variances},\\[2pt]
\bigl(\bm{\Sigma}_{ij}\,N/\operatorname{tr}\bm{\Sigma}_{ij}\bigr)^{-1} & \text{unit trace}.
\end{cases}
\label{eq:norm-variants}
\end{equation}
%
The first and fourth remove the \emph{size} of a cluster while retaining its
shape, so that clusters of different populations compete on equal terms; for an
isotropic cluster the first reduces to the Euclidean distance.

\subsection{Local models}
\label{subsec:ecr-networks}

Each cluster carries a Gaussian radial-basis-function network mapping a state
to a normalised perturbation. With standardised inputs
$\tilde{\mathbf{z}}=(\mathbf{z}-\bar{\mathbf{z}})\oslash\mathbf{s}$, centres
$\mathbf{c}_k$ and widths $s_k$,
%
\begin{equation}
\phi_k(\mathbf{z})=\exp\!\left(-\frac{\|\tilde{\mathbf{z}}-\mathbf{c}_k\|^{2}}{2 s_k^{2}}\right),
\qquad k=1,\dots,n_c,
\label{eq:rbf-basis}
\end{equation}
%
\begin{equation}
\widehat{\Delta\mathbf{p}}=\mathbf{W}^{\mathsf T}\bigl[\phi_1(\mathbf{z})\;\cdots\;\phi_{n_c}(\mathbf{z})\;\;1\bigr]^{\mathsf T},
\qquad
\widehat{\mathbf{p}}=\mathbf{p}_{\mathrm{nom}}+\widehat{\Delta\mathbf{p}}\odot\delta\mathbf{p}_{\max}.
\label{eq:rbf-output}
\end{equation}
%
Centres are placed by $k$-means on the cluster data, widths from the mean
distance to the nearest centres, and the output weights follow in closed form
from regularised least squares,
%
\begin{equation}
\mathbf{W}=\bigl(\bm{\Phi}^{\mathsf T}\bm{\Phi}+\lambda\mathbf{I}\bigr)^{-1}\bm{\Phi}^{\mathsf T}\mathbf{Y},
\qquad
\bm{\Phi}_{nk}=\phi_k(\mathbf{z}_n),\quad
\mathbf{Y}_{n:}=\Delta\mathbf{p}_n^{\mathsf T},
\label{eq:rbf-weights}
\end{equation}
%
which is why ECR training costs seconds rather than the iterative effort of a
back-propagation network. Because the target of the regression is a normalised
perturbation, admissibility is the box test of
\eqref{eq:normalised-perturbation}, and the membership oracle used in
\eqref{eq:index-Ti} is simply
%
\begin{equation}
\mathcal{M}_i(\mathbf{z})=
\begin{cases}
1, & \text{if } \bigl\|\widehat{\Delta\mathbf{p}}^{\,(i j^{*})}(\mathbf{z})\bigr\|_{\infty}\le 1
\;\text{ for the cluster } j^{*}=\arg\min_j \mathrm{NMD}_{ij}(\mathbf{z}),\\
0, & \text{otherwise,}
\end{cases}
\label{eq:membership-oracle}
\end{equation}
%
optionally restricted to clusters for which
$\mathrm{NMD}_{ij}(\mathbf{z})\le\gamma\,q_{95}(C_{ij})$, where $q_{95}$ is the
$95$th percentile of the NMD of the cluster's own training data and $\gamma$ a
gate factor. The gate exists because \eqref{eq:nmd} always returns a nearest
cluster, including for states nowhere near any recorded data, for which the
network extrapolates.

\subsection{On-line control law}
\label{subsec:ecr-online}

At each crossing the controller measures $\mathbf{z}_n$, identifies a region,
and applies the perturbation that region prescribes. The three variants differ
only in how the region is identified.

\paragraph{ECR-I.} One network is trained over the whole activation region and
evaluated directly; the current region is never identified explicitly.

\paragraph{ECR-II.} The state is offered to the networks in ascending region
order and the first admissible answer is accepted:
%
\begin{equation}
i^{*}=\min\bigl\{i\in\{0,\dots,K\}:\|\widehat{\Delta\mathbf{p}}^{\,(i)}(\mathbf{z}_n)\|_{\infty}\le1\bigr\}.
\label{eq:ecr2-law}
\end{equation}

\paragraph{ECR-III.} The region follows analytically from \eqref{eq:nmd}, the
search running over the clusters of \emph{all} regions at once:
%
\begin{equation}
(i^{*},j^{*})=\arg\min_{i,j}\;\mathrm{NMD}_{ij}(\mathbf{z}_n).
\label{eq:ecr3-law}
\end{equation}

In every variant the applied parameter is
%
\begin{equation}
\mathbf{p}_n=
\begin{cases}
\mathbf{p}_{\mathrm{nom}}+\widehat{\Delta\mathbf{p}}^{\,(i^{*}j^{*})}(\mathbf{z}_n)\odot\delta\mathbf{p}_{\max},
& \text{if such an } (i^{*},j^{*}) \text{ exists},\\[2pt]
\mathbf{p}_{\mathrm{nom}}, & \text{otherwise (no targeting)} ,
\end{cases}
\label{eq:control-law}
\end{equation}
%
so the perturbation is admissible by construction and the plant runs at nominal
parameters whenever the state is unrecognised.

\subsection{Performance measures}
\label{subsec:ecr-metrics}

For an initial condition $\mathbf{z}_0$ the \emph{reaching time} is the number
of map steps needed to enter the target region,
%
\begin{equation}
n_r(\mathbf{z}_0)=\min\bigl\{n\ge0:\|\mathbf{z}_n-\mathbf{z}^{*}\|<\delta\bigr\},
\label{eq:reaching-time}
\end{equation}
%
and the reported figure of merit is its mean over a fixed ensemble
$\{\mathbf{z}_0^{(m)}\}_{m=1}^{M_{\mathrm{ic}}}$ of initial conditions drawn
from the attractor, censored at a horizon $n_{\max}$,
%
\begin{equation}
\bar{n}_r=\frac{1}{M_{\mathrm{ic}}}\sum_{m=1}^{M_{\mathrm{ic}}}
\min\bigl\{n_r(\mathbf{z}_0^{(m)}),\,n_{\max}\bigr\} .
\label{eq:mean-reaching-time}
\end{equation}
%
Two further quantities separate targeting from stabilisation: the
\emph{capture rate} $\Pr\{n_r<n_{\max}\}$, and the \emph{retention}
%
\begin{equation}
\eta=\frac{1}{n_h}\sum_{n=n_r+1}^{n_r+n_h}\mathbb{1}\bigl[\|\mathbf{z}_n-\mathbf{z}^{*}\|<\delta\bigr],
\label{eq:retention}
\end{equation}
%
the fraction of the $n_h$ steps following capture during which the local
controller holds the state inside the target region. Without targeting,
$\bar{n}_r$ measures the natural recurrence time of the section trajectory to
$\mathcal{B}_\delta$; for the present configuration only $6\%$ of crossings
fall inside $\mathcal{B}_\delta$, and $\bar{n}_r\approx15$ steps.

\subsection{Algorithm}
\label{subsec:ecr-algorithm}

\noindent\textbf{Off-line (training).}
\begin{enumerate}
\item Collect $\mathcal{D}$ by applying $\mathbf{p}_n\sim\mathcal{U}(\Pi)$ at
every crossing, \eqref{eq:dataset}.
\item Label $\mathcal{I}_0$ by \eqref{eq:index-T0}; cluster $\mathcal{Z}_0$
with radius $r_0$ from \eqref{eq:distance-histogram}; compute
$(\bm{\mu}_{0j},\underline{\mathbf{N}}_{0j})$ by
\eqref{eq:cluster-stats}--\eqref{eq:norm-variants}; fit $NN_{0j}$ by
\eqref{eq:rbf-weights}.
\item For $i=1,\dots,K$: label $\mathcal{I}_i$ by \eqref{eq:index-Ti} using the
oracle \eqref{eq:membership-oracle} of level $i-1$; cluster; fit $NN_{ij}$.
Stop when $|\mathcal{I}_i|$ falls below a minimum count.
\end{enumerate}

\noindent\textbf{On-line (targeting and control).}
\begin{enumerate}
\item Measure $\mathbf{z}_n$ at the crossing.
\item Identify the region by \eqref{eq:ecr2-law} or \eqref{eq:ecr3-law}.
\item Apply $\mathbf{p}_n$ from \eqref{eq:control-law} and hold it until the
next crossing.
\end{enumerate}

\subsection{Remarks on the underdetermined elements}
\label{subsec:ecr-remarks}

\begin{remark}[Choices not fixed by the source formulation]
The normalisation in \eqref{eq:norm-variants}, the gate factor $\gamma$, the
tie-breaking rule in \eqref{eq:ecr3-law}, and the selection of the training
pairs for $NN_{0j}$ are all left open in the original description. They are not
cosmetic: on the Lorenz section the size-removing normalisations reach the
target in $4.3$--$4.4$ steps with retention $0.96$, whereas plain Mahalanobis
or normalised-variance readings require $5.2$ steps and hold the orbit only
$0.61$--$0.69$ of the time, with markedly larger run-to-run spread.
\end{remark}

\begin{remark}[Delay coordinates]
\label{rem:delay}
Under \eqref{eq:delay-embedding} the components $s(t_n-kT)$, $k\ge1$, were
generated while the \emph{previous} parameter vector was applied. The observed
transition is therefore a function of $(\mathbf{z}_n,\mathbf{p}_n,\mathbf{p}_{n-1})$
rather than of $(\mathbf{z}_n,\mathbf{p}_n)$ alone, and \eqref{eq:map} holds
only approximately. Coarse targeting tolerates this ambiguity --- the reaching
time still improves substantially --- but local stabilisation does not, the
retention \eqref{eq:retention} falling from $\approx0.96$ in real coordinates
to $\approx0.20$. The standard remedy is to augment the observed state with the
previously applied perturbation, $[\mathbf{z}_n^{\mathsf T}\;\;\Delta\mathbf{p}_{n-1}^{\mathsf T}]^{\mathsf T}$.
\end{remark}

\begin{remark}[Deterministic deadlock]
Because \eqref{eq:control-law} is a deterministic state feedback, a trajectory
can be locked into a \emph{controlled} periodic orbit that never enters
$\mathcal{B}_\delta$, inflating \eqref{eq:mean-reaching-time} through censored
samples rather than through slow targeting. Applying $\mathbf{p}_{\mathrm{nom}}$
for a single step after a fixed number of unsuccessful controlled steps, or the
presence of measurement noise, removes the deadlock.
\end{remark}
